\section{The same Boolean resource for tanh and sigmoid}
\label{sec:saturating-Boolean}
Bounded saturating activations give an even simpler digital implementation. In this section no small-intensity condition on the scalar prior is required.

\begin{proposition}[Saturating Boolean gates]
\label{prop:saturating-Boolean}
The Boolean expected-risk and online sign-prediction conclusions of Section~\ref{sec:digital-circuits} also hold for tanh and logistic sigmoid hidden activations. A layered fan-in-two Boolean circuit of depth $h$ needs at most $h$ hidden layers and one additional constant wire per layer. Padding uses saturated COPY gates. The learner width must contain this implementation. Coefficients and row fan-in are bounded by numerical constants, and the resource is again
\[
 J_*\log N+E_{1,*}\log(d+1)
       +C(S+H+1)\log\{C(n+1)(S+H+1)\}.
\]
Here the counts are those of the saturating implementation, so they need not equal the ReLU counts. The observation likelihood and prediction use a final clipping link. The exact sign-prediction conclusion instead uses the all-input raw margin convention already specified.
\end{proposition}
\begin{proof}
For tanh encode true and false by values in $[0.9,1]$ and $[-1,-0.9]$. Carry a constant in $[0.9,1]$. The gate preactivations are
\[
 4(a+b-c)\quad\hbox{for AND},\qquad -4a\quad\hbox{for NOT},
 \qquad4a\quad\hbox{for COPY},\qquad4c\quad\hbox{for the constant}.
\]
For AND, the expression $a+b-c$ is at least $0.8$ when both inputs are true and at most $-0.8$ otherwise. The other margins are larger. Any added preactivation perturbation of magnitude at most $0.1$ therefore preserves the encoding intervals after tanh. The first layer uses the exact Boolean input signs and constant coordinate, so it starts the induction. The readout $2a$ has the correct sign and magnitude at least $1.8$ before a readout perturbation; clipping it gives the exact Boolean target whenever that perturbation is at most $0.1$.

For the logistic sigmoid $\sigma(t)=(1+e^{-t})^{-1}$, encode false in $[0,0.1]$ and true in $[0.9,1]$, with a constant in $[0.9,1]$. Use preactivations
\[
 20(a+b-1.5c),\quad20(0.5c-a),\quad20(a-0.5c),\quad20c
\]
for AND, NOT, COPY and the constant. The corresponding preactivation signs have absolute margins at least five before a perturbation of size $0.1$, and the output remains in its prescribed encoding interval. Initially replace each input bit by $(1+x_i)/2$ inside the first gate's affine expression, so no extra hidden layer is needed. The readout $4a-2c$ has magnitude at least $1.4$ with the correct sign. Again clipping is exactly correct on the good event. All coefficients and numbers of used input coordinates per row are bounded by numerical constants.

For either activation all selected and native hidden features are bounded in absolute value by one. In the same comparison construction, every selected-row or readout residual therefore has expected uniform magnitude at most $C\epsilon$, at every depth, without a condition on the native row moment. A union bound and Markov's inequality show that all these residuals are at most $0.1$ with probability at least $1-C(V_*+1)\epsilon$. On this event the clipped output is exactly the target everywhere on the Boolean cube. Taking $\epsilon=ct/(V_*+1)$ gives expected uniform clipped-output error at most $t$ and entropy \eqref{eq:digital-entropy}, with the new resource counts. The margin lower bound uses a fixed positive success probability. Apply the same statistical and online arguments as before.
\end{proof}
